\section{Five consecutive regions and their original bonds}

\begin{definition}[Internal bonds and regional sums]
    \label{def:five_region_bond_partition}
    \uses{def:chain_model}
    For a chain of $L$ sites put $\mathcal B_L=\{i\in\N:i+1<L\}$ and,
    for arbitrary $l,h\in\N$, set
    \begin{align}
        R(l,h) & =\{v<L:l\le v<h\},\notag              \\
        I(l,h) & =\{i\in\mathcal B_L:l\le i,\ i+1<h\}.
    \end{align}
    Given $k_1\le k_2\le k_3\le k_4$, put $k_0=0$, $k_5=L$ and
    $I_j=I(k_j,k_{j+1})$, $A_j=\sum_{i\in I_j}h_i$ for $0\le j<5$.
    Set $J=\{i\in\mathcal B_L:i+1\in\{k_1,k_2,k_3,k_4\}\}$ and
    $Y=\sum_{i\in J}h_i$. All site sets are intersected with the finite chain.
\end{definition}

\begin{theorem}[Support of an internal bond]
    \label{thm:internal_bond_region_subset}
    \lean{TNLean.MPS.AreaLaw.internal_bond_support_subset}
    \uses{def:five_region_bond_partition}
    If $i\in I(l,h)$, then $\{i,i+1\}\subseteq R(l,h)$.
\end{theorem}
\begin{proof}
    The inequalities $l\le i$ and $i+1<h$ hold by definition; both sites
    lie in the chain because $i+1<L$.
\end{proof}

\begin{theorem}[The number of internal bonds]
    \label{thm:internal_bond_card_bound}
    \lean{TNLean.MPS.AreaLaw.card_internal_bonds_le}
    \uses{def:five_region_bond_partition}
    For arbitrary $L,l,h\in\N$, one has
    $|I(l,h)|\le\max\{h-l-1,0\}$, where subtraction inside the maximum
    is integer subtraction. Empty and out-of-range intervals are included.
\end{theorem}
\begin{proof}
    If $h\le l+1$ there are no internal bonds. Otherwise every bond start
    lies between $l$ and $h-2$, giving at most $h-l-1$ choices.
\end{proof}

\begin{theorem}[Commutation of consecutive internal sums]
    \label{thm:consecutive_internal_sums_commute}
    \lean{TNLean.MPS.AreaLaw.pairwise_commute_sum_internal_consecutive_regions}
    \uses{def:five_region_bond_partition}
    Let $k_0\le\cdots\le k_m$ be natural numbers. If every $h_i$ is
    supported on $\{i,i+1\}$, the operators
    $\sum_{i\in I(k_j,k_{j+1})}h_i$, $0\le j<m$, commute pairwise.
    No positivity or Hermiticity assumption is needed.
\end{theorem}
\begin{proof}
    \uses{thm:internal_bond_region_subset}
    Each sum is supported on $R(k_j,k_{j+1})$. These regions are pairwise
    disjoint, so their supported operators commute.
\end{proof}

\begin{theorem}[The additive five-region decomposition]
    \label{thm:five_region_additive_partition}
    \lean{Finset.sum_five_internal_regions_add_crossings}
    \uses{def:five_region_bond_partition}
    For $k_1\le k_2\le k_3\le k_4$ and arbitrary values $h_i$ in an
    additive commutative monoid,
    \begin{align}
        \sum_{i\in\mathcal B_L}h_i=\sum_{j=0}^4A_j+Y.
    \end{align}
    No upper bound on the cut positions is required for this identity.
\end{theorem}
\begin{proof}
    A bond either has its right endpoint at one of the four cuts or is
    internal to exactly one of the five regions. Partition the finite sum
    according to these alternatives. Repeated cuts contribute only once
    to the crossing set.
\end{proof}

\begin{theorem}[Original regional support and order budgets]
    \label{thm:five_region_original_bond_bounds}
    \lean{TNLean.MPS.AreaLaw.five_region_original_bond_data}
    \uses{def:five_region_bond_partition}
    Suppose $k_4\le L$, $0\le h_i\le\Id$, and $h_i$ is supported on
    $\{i,i+1\}$. Then $(k_j)_{j=0}^5$ is monotone, the regions
    $R_j=R(k_j,k_{j+1})$ are pairwise disjoint, and $A_j$ is supported
    on $R_j$. In addition,
    \begin{align}
        0\le A_j & \le |I_j|\Id, & 0\le Y     & \le |J|\Id,\notag   \\
        |J|      & \le4,         & \sum_i h_i & =\sum_{j=0}^4A_j+Y.
    \end{align}
\end{theorem}
\begin{proof}
    \uses{thm:internal_bond_region_subset,thm:five_region_additive_partition}
    Support follows term by term, and positive unit bounds add.
    Each crossing bond is determined by its right endpoint, so there are
    at most four. Apply the additive partition to the original operators.
\end{proof}

\begin{theorem}[The seven-scale clipping budget]
    \label{thm:five_region_seven_ell_budget}
    \lean{Real.five_region_final_clipping_budget_le}
    Let $\ell\ge1$, let $c_0,\ldots,c_4,C$ be natural numbers with
    $c_1,c_3\le2\ell-2$ and $C\le4$, and put $K=\{0,2,4\}$.
    Then $\sum_{j\in K}\ell+\sum_{j\notin K}c_j+C\le7\ell$.
\end{theorem}
\begin{proof}
    The left side is $3\ell+c_1+c_3+C$, at most
    $3\ell+2(2\ell-2)+4=7\ell$.
\end{proof}
